% -----------------------------------------------------------------------------
%  00_abstract.tex
%  Every numeric value here is a macro from macros.tex (authority: NUMBERS.md).
% -----------------------------------------------------------------------------

We set out to reproduce a preprint's claim that a masked diffusion language model
extracts structured records from logs \paperClaimedSpeedup$\times$ faster than a
quality-matched autoregressive (AR) baseline. The reproduction collapsed, and it
collapsed on our side of the comparison: our AR baseline carried two independent
defects, each flattering diffusion. Qwen2.5 ships a repetition penalty in its
generation configuration that HuggingFace \texttt{generate()} applies even under
greedy decoding --- and structured extraction copies spans verbatim out of the
input, precisely what such a penalty suppresses. Separately, our AR latencies
came from a naive \texttt{generate()} loop while the preprint serves its own AR
baselines on a paged-attention stack. Neither defect could touch the diffusion
side. Corrected, the headline reverses: a properly served AR-1.5B is
\arSpeedupOverMDM$\times$ faster than LLaDA-8B and \parityGap\ Field-F1 better.

We then asked what \emph{is} true, and fitted two cost laws on a shared serving
stack: AR latency is linear in output length (\arSlope~ms/token), diffusion
latency stepwise in canvases of \dgCanvas\ tokens (\dgStep~ms per additional
canvas, \dgMarginal~ms/token at the margin). They are two orthogonal walls ---
diffusion is flat in length and walls in arrival rate at
$\lambda \approx \dgKneeLo$--\dgKneeHi~ev/s; AR is flat in rate to
$\lambda \approx \arKneeLo$--\arKneeHi~ev/s and walls in length --- and both
follow from one mechanism: bidirectional attention invalidates the KV-cache
reuse that makes AR batching cheap, leaving diffusion compute-bound at batch~1
with no headroom for batching to convert (multiplier
\arMultOneTwoEight$\times$ for AR against \dgMultOneTwoEight$\times$). We
confirm the cap is architectural, not an engine artifact, by reproducing it
(\lladaMultDinfer$\times$) on a second, diffusion-native runtime. Diffusion's
advantage is therefore not latency but invariance to output length in
\emph{both} cost and quality, realizable only where batching is unavailable. We
test six regimes the laws predict: four confirm, two fail --- including our own
prediction of perturbation robustness, refuted in \perturbCategories\ of
\perturbCategories\ categories.
